% TOP_PROBLEM: {"id": 2100710, "title": "Open Problems in Integrable Systems — Quantum integrability for polynomial in momenta integrals", "category_id": 11, "category": {"id": 11, "name": "dynamical_systems", "display_name": "Dynamical Systems", "description": "Problems about long-term behavior of deterministic systems, Hamiltonian dynamics, and periodic orbits.", "slug": "dynamical-systems", "order_index": 11, "created_at": "2026-07-31T15:26:25.671Z"}, "status": "open"}
% FIRST_POSTED: null
% ENTRY_KIND: statement_only
% Imported from: batch09.tex; UnsolvedMath ID 2100710
\documentclass[11pt]{article}
\usepackage[margin=25mm]{geometry}
\usepackage{fontspec}
\setmainfont{Latin Modern Roman}
\usepackage{amsmath,amssymb,mathtools}
\usepackage{unicode-math}
\setmathfont{Latin Modern Math}
\usepackage{xurl,hyperref}
\hypersetup{colorlinks=true,urlcolor=blue}
\setcounter{secnumdepth}{0}
\setlength{\parindent}{0pt}
\setlength{\parskip}{5pt}
\setlength{\emergencystretch}{3em}
\newcommand{\sref}[2]{\hyperlink{source:#1:#2}{[S#2]}}
\newcommand{\cataloguescope}[1]{\paragraph{Recorded target.}#1}
\begin{document}
% UnsolvedMath ID 2100710; source catalogue code AMR-020-0710
\setcounter{section}{2100709}
\section{Open Problems in Integrable Systems — Quantum integrability for polynomial in momenta integrals}\label{2100710}
{\small\sffamily UnsolvedMath ID 2100710 · Dynamical Systems}
\subsection{Definitions and mathematical statement}

\paragraph{Shared notation.}$\mathbb N=\{1,2,\ldots\}$, and $\mathbb Z,\mathbb Q,\mathbb R,\mathbb C$ denote the integers, rationals, reals and complex numbers. Any other conventions are specified locally.


Let $(M,g)$ be a smooth Riemannian manifold, with geodesic Hamiltonian $K(x,p)=\tfrac12g^{ij}(x)p_ip_j$ on $T^*M$ and its canonical Poisson bracket. A polynomial geodesic first integral is a smooth function polynomial in the cotangent variables, satisfying $\{K,F\}=0$. Its homogeneous pieces are themselves first integrals. The Laplace--Beltrami operator is $\Delta_g=g^{ij}\nabla_i\nabla_j$. A homogeneous integral of degree $d$ is quantizable if there is a differential operator $\widehat F$ of order $d$, with principal symbol $F$ and arbitrary allowed lower-order corrections, such that
\[[\Delta_g,\widehat F]=\Delta_g\widehat F-\widehat F\Delta_g=0.\]
Quantization of a general polynomial integral means quantization of its homogeneous components; self-adjointness is not imposed. The principal symbol is the top-order homogeneous polynomial in the differential operator's cotangent variables.
\paragraph{Statement recorded in the supplied source.}
Find necessary and/or sufficient conditions on $g$ ensuring that every polynomial first integral of its geodesic flow is quantizable in this sense, without restricting the integral degree or the dimension of $M$. \sref{2100710}{2}
\subsection{Short English statement}
Which metrics allow every polynomial integral of geodesic motion to become an exactly commuting differential operator with the same principal symbol?
\subsection{Sources}
\begin{description}
\item[\hypertarget{source:2100710:1}{S1}] ULAM.AI and the UnsolvedMath Contributors, UnsolvedMath: A Curated Collection of Open Mathematics Problems, record 2100710 (AMR-020-0710). CC BY 4.0. \url{https://huggingface.co/datasets/ulamai/UnsolvedMath}.
\item[\hypertarget{source:2100710:2}{S2}] Alexey Bolsinov, Vladimir S. Matveev, Eva Miranda and Serge Tabachnikov, Open Problems, Questions, and Challenges in Finite-Dimensional Integrable Systems (2018), Problem 7.11 and its complete preceding definitions. Original question and full discussion. \url{https://arxiv.org/abs/1804.03737}.

\end{description}
\label{end:2100710}
\end{document}
